\chapter{Requirements, Development Lifecycle and System Architecture}
\label{chap:requirements-architecture}

This chapter defines the problem addressed by the framework, its scope, its
requirements, the development lifecycle followed during the project, and the
architecture adopted by the final implementation. The framework accepts a YAML
configuration, validates it, and produces a standard Kathar\'a laboratory. It
also provides lifecycle scripts through which a generated laboratory can be
started, stopped, inspected, and removed. The detailed YAML model and the
validation rules are presented in Chapter~\ref{chap:configuration-validation},
while the translation of a valid configuration into \texttt{lab.conf}, startup
files, FRRouting files, and node resources is discussed in the following
chapter. The purpose of this chapter is therefore not only to state what the
framework does, but also to explain why its main components were introduced,
how problems encountered during development changed the requirements, and how
the resulting responsibilities cooperate as a single system.

\section{Problem, Scope, and Requirements}
\label{sec:problem-scope-requirements}

Kathar\'a provides a flexible representation of network laboratories
\cite{kathara_website,kathara_manpage}, but the configuration of a non-trivial
topology is normally distributed across several artefacts. A node interface is declared in \texttt{lab.conf}; its IP address is
assigned in a startup file; the corresponding prefix may appear again in an
FRRouting configuration; and scripts or other auxiliary resources may be
stored in the directory associated with the node. Each individual file is
readable, but the consistency of the complete laboratory must usually be
maintained manually.

This becomes progressively more difficult as the topology grows. A change to a
single address may require coordinated changes to a gateway, a routing
neighbour, an advertised network, a static route, or another reference used by
a node. The individual files may remain syntactically valid even when the
laboratory is no longer semantically coherent, and the resulting error may
only become visible after the emulated environment has started. Repetition is
a related problem: hosts repeatedly require interface configuration and a
default route, while routers repeatedly require interface addressing, an
appropriate container image, and routing configuration. Rewriting the same
information in multiple files increases the probability of inconsistent node
names, IP addresses, and interface indexes.

A further limitation of the standard file-based workflow is the absence of an
explicit validation boundary. Kathar\'a configuration files are executable
artefacts rather than a high-level schema that can be analysed before output
is written. The proposed framework introduces such a boundary by treating a
single declarative YAML document as the source configuration, validating that
document before generation, and only then producing the corresponding
Kathar\'a files. The general problem addressed by the work can therefore be
stated as follows:

\begin{quote}
How can a network laboratory be described once, checked for structural and
semantic consistency, translated into standard Kathar\'a artefacts, and
managed through a repeatable lifecycle?
\end{quote}

The framework is intentionally limited in scope. It does not replace Kathar\'a,
Docker, FRRouting, or Wireshark; instead, it acts as an automation layer that
prepares input for those tools. The implemented scope includes YAML-based
descriptions of laboratories, nodes, interfaces, addressing, gateways, static
routes, custom commands, and routing protocols; structural and semantic
validation before generation; generation of \texttt{lab.conf} and one startup
file per node; generation of FRRouting files for OSPF, RIP, and BGP; optional
import of directories associated with declared nodes; optional integration of
a real-time Wireshark node; explicit handling of pre-existing output
directories; shell wrappers for starting, stopping, and deleting laboratories;
and support for multiple configurations and generated laboratories. The
experimental components contained in the repository use these generated
laboratories, but they are not part of the core configuration architecture
discussed here; their workloads, measurements, and evaluation are addressed in
later chapters.

The primary user is a student, researcher, or network experimenter who already
understands basic Linux networking and needs to create repeatable Kathar\'a
laboratories. The framework does not hide networking concepts: the user still
chooses addresses, prefixes, gateways, routing protocols, and neighbours. Its
purpose is to avoid manually distributing the same design choices among many
files. A secondary stakeholder is the reviewer of an experiment. Generated
files must remain readable so that the relationship between the YAML source
and the resulting Kathar\'a configuration can be inspected. This requirement
motivates the generation of ordinary text files rather than the use of a
private runtime data store. The expected execution environment is a Linux
workstation with Python, PyYAML, Docker, and Kathar\'a installed
\cite{pyyaml_docs,docker_overview,kathara_website}. FRRouting is
provided by the appropriate Kathar\'a image, while Wireshark integration also
requires the selected container image and support for the exposed interface.

The resulting workflow can be summarised through a small set of principal use
cases. The user may select a configuration by supplying either an explicit
path or a configuration name resolved inside a configured search directory.
The framework then loads the document with PyYAML and requires the top-level
value to be a mapping. The user may validate the configuration without
creating or modifying a laboratory: structural validation checks mandatory
fields, while semantic validation checks relationships among addresses,
collision domains, gateways, roles, and routing declarations, producing a
report containing errors, warnings, and summary lines. After successful
validation, the user may generate the laboratory under the selected output
directory; the framework writes \texttt{lab.conf}, one startup file per node,
and the required FRRouting files, while existing output is handled according
to the selected \cliopt{force} or \cliopt{clean} behaviour. If directory
import is enabled, a directory located beside the YAML source is copied only
when its name matches a declared node; unmatched directories are reported as
warnings and skipped. Wireshark can be disabled, selected interactively,
attached directly to one or more collision domains, or associated with the
networks connected to a specific node. Finally, the generated laboratory can
be started, inspected, cleaned, and removed without deleting the original YAML
configuration.

These use cases are formalised by the following functional requirements.

\begin{longtable}{p{0.12\textwidth}p{0.78\textwidth}}
\caption{Functional requirements of the framework.}
\label{tab:functional-requirements}\\
\toprule
\textbf{ID} & \textbf{Requirement} \\
\midrule
\endfirsthead
\toprule
\textbf{ID} & \textbf{Requirement} \\
\midrule
\endhead
FR1 & The framework shall accept either an explicit YAML path or a configuration name resolved inside a configurable directory. \\
FR2 & The framework shall load YAML through a safe parser and reject a top-level value that is not a mapping. \\
FR3 & The framework shall require a laboratory name and a non-empty node mapping. \\
FR4 & The framework shall validate identifiers, node types, interfaces, collision domains, addresses, gateways, and routing declarations. \\
FR5 & The framework shall distinguish generation-blocking errors from non-blocking warnings. \\
FR6 & The framework shall provide a validation-only mode that does not create or modify a laboratory. \\
FR7 & The framework shall generate a Kathar\'a \texttt{lab.conf} file from node-interface declarations. \\
FR8 & The framework shall select a default container image according to node capabilities and permit an explicit image override. \\
FR9 & The framework shall generate one startup file for every declared node. \\
FR10 & Startup generation shall support interface addresses, a default gateway, static routes, FRR startup, and custom commands. \\
FR11 & The framework shall generate FRRouting daemon and configuration files for enabled FRR devices. \\
FR12 & The framework shall support OSPF, RIP, and BGP configuration forms represented by the YAML model. \\
FR13 & The framework shall optionally copy node directories located beside the source YAML file. \\
FR14 & The framework shall optionally attach a Wireshark node to selected collision domains. \\
FR15 & The framework shall avoid silently overwriting an existing laboratory unless the user selects an explicit overwrite or clean-regeneration option. \\
FR16 & The framework shall support a configurable output directory containing multiple generated laboratories. \\
FR17 & The project shall provide commands or wrappers to start, stop, inspect, and delete generated laboratories. \\
\bottomrule
\end{longtable}

In addition to these functions, the framework is constrained by a set of
non-functional requirements that influence both the architecture and the
implementation style.

\begin{longtable}{p{0.22\textwidth}p{0.68\textwidth}}
\caption{Non-functional requirements of the framework.}
\label{tab:non-functional-requirements}\\
\toprule
\textbf{Property} & \textbf{Requirement} \\
\midrule
\endfirsthead
\toprule
\textbf{Property} & \textbf{Requirement} \\
\midrule
\endhead
Reproducibility & The same maintained source configuration should regenerate the same logical laboratory structure. \\
Transparency & Generated files must remain ordinary, readable Kathar\'a and FRRouting artefacts. \\
Safety of file operations & Existing laboratory directories must not be removed or overwritten implicitly. \\
Modularity & Loading, validation, generation, monitoring selection, and lifecycle operations should have separate responsibilities. \\
Extensibility & New configuration properties or protocol forms should be addable without replacing the entire generation pipeline. \\
Usability & Common operations should be available through a command-line interface with explicit options and informative messages. \\
Traceability & Validation messages should identify the affected node, interface, domain, neighbour, or field. \\
Portability & The project should operate in a standard Linux, Python, Docker, and Kathar\'a environment. \\
Maintainability & Persistent modifications should be made in source YAML or resource directories rather than only in generated output. \\
Compatibility & The generated laboratory should remain operable with standard Kathar\'a commands. \\
\bottomrule
\end{longtable}

\subsection*{Requirements Traceability and Verification}
\phantomsection
\label{sec:functional-requirements-traceability}

Table~\ref{tab:functional-requirements-traceability} links each functional
requirement to the implementation behaviour and to the evidence presented in
the thesis. ``Verified'' denotes that the behaviour is demonstrated by the
described validation or generation workflow and, where applicable, by
execution of generated laboratories. It does not imply exhaustive testing of
every possible input combination.

\begin{longtable}{p{0.08\textwidth}p{0.64\textwidth}p{0.16\textwidth}}
\caption{Traceability matrix for functional requirements FR1--FR17.}
\label{tab:functional-requirements-traceability}\\
\toprule
\textbf{ID} & \textbf{Implementation and verification evidence} & \textbf{Assessment} \\
\midrule
\endfirsthead
\toprule
\textbf{ID} & \textbf{Implementation and verification evidence} & \textbf{Assessment} \\
\midrule
\endhead
FR1 & Path and configuration-name resolution are described in Section~\ref{sec:configuration-model-purpose} and exercised by the generation workflow. & Verified \\
FR2 & Safe YAML loading and rejection of a non-mapping root are part of the structural-validation boundary in Section~\ref{sec:structural-validation}. & Verified \\
FR3 & Mandatory \texttt{lab\_name} and non-empty \texttt{nodes} checks are documented with the structural rules in Section~\ref{sec:structural-validation}. & Verified \\
FR4 & Identifier, role, interface, addressing, gateway, collision-domain, and routing checks are described in Sections~\ref{sec:structural-validation} and~\ref{sec:validation-workflow}. & Verified \\
FR5 & The error/warning policy and its effect on generation are reported in Table~\ref{tab:error-warning-policy}. & Verified \\
FR6 & Validation-only execution terminates before output generation, as described in Section~\ref{sec:validation-workflow}; the large invalid stress case also produced no files in Section~\ref{sec:stress-methodology}. & Verified \\
FR7 & Translation from interface declarations to \texttt{lab.conf} is shown in Section~\ref{sec:generated-laboratory-structure} and Table~\ref{tab:yaml-to-generated}. & Verified \\
FR8 & Capability-based default-image selection and explicit overrides are documented in Sections~\ref{sec:laboratory-generation-process} and~\ref{sec:generated-laboratory-structure}. & Verified \\
FR9 & Per-node startup-file generation is described and illustrated in Section~\ref{sec:generated-laboratory-structure}. & Verified \\
FR10 & Generated address, gateway, route, FRR-startup, and custom-command entries are mapped in Table~\ref{tab:yaml-to-generated}. & Verified \\
FR11 & Generation of \texttt{daemons} and \texttt{frr.conf} for FRR-enabled nodes is described in Section~\ref{sec:generated-laboratory-structure}. & Verified \\
FR12 & The supported OSPF, RIP, and BGP forms and their validation are documented in Sections~\ref{sec:configuration-model-purpose} and~\ref{sec:structural-validation}; routed laboratories provide runtime use evidence. & Verified \\
FR13 & Selective import of directories whose names match declared nodes is described in Sections~\ref{sec:laboratory-generation-process} and~\ref{sec:generated-laboratory-structure}. & Verified \\
FR14 & Wireshark selection, generated attachment, and a runtime capture are presented in Section~\ref{sec:generated-laboratory-structure} and Figure~\ref{fig:wireshark-runtime-capture}. & Verified \\
FR15 & Refusal, explicit overwrite, and clean-regeneration behaviours are documented in Section~\ref{sec:laboratory-generation-process}. & Verified \\
FR16 & Output-root selection and coexistence of named generated laboratories are described in Sections~\ref{sec:laboratory-generation-process} and~\ref{sec:laboratory-lifecycle}. & Verified \\
FR17 & Start, stop, status, connection, clean, and deletion operations and their effects are reported in Section~\ref{sec:laboratory-lifecycle} and Table~\ref{tab:lifecycle-effects}. & Verified \\
\bottomrule
\end{longtable}

The matrix makes two boundaries explicit. First, FR12 is verified only for the
protocol forms represented by the current YAML model, not for every command
offered by FRRouting. Second, runtime execution complements but does not replace
static verification: opaque custom commands and complete manual FRR fragments
remain outside the semantics that the validator can prove.

\section{Software Development Lifecycle and Design Evolution}
\label{sec:software-development-lifecycle}

The final architecture should not be interpreted as the implementation of a
complete specification that existed before coding began. Development followed
an iterative and incremental lifecycle in which a small working solution was
extended after each cycle of generation, execution, observation, and
correction. This approach was appropriate for the project because several
requirements only became evident when real Kathar\'a laboratories were produced
and executed. In particular, failures that were invisible at the level of YAML
syntax or generated text files often became apparent only when routes,
services, or measurement scripts were exercised at runtime.

The initial requirement was deliberately narrow. Manually preparing a
Kathar\'a laboratory required the same design information to be distributed
among \texttt{lab.conf}, node startup files, FRRouting configuration, and
node-specific directories. One possible approach would have been to automate
edits to these files individually, but that would have preserved the very
duplication that motivated the project. A single declarative source was
therefore preferred. YAML was selected because its hierarchical structure maps
naturally to laboratories, nodes, interfaces, and nested routing properties,
while remaining readable and directly editable. The first useful prototype
could consequently describe a topology at a higher level and derive ordinary
Kathar\'a files from it.

Generation alone soon proved insufficient. A YAML document could be parsed and
translated while still describing an inconsistent network: for example, an
address could be duplicated, a default gateway could lie outside every local
subnet, or a routing relationship could refer to the wrong peer. Allowing such
configurations to reach the runtime moved deterministic errors to a much later
and less convenient debugging stage. Validation therefore became an explicit
architectural responsibility. Structural validation was introduced to protect
the implementation from malformed document shapes, while semantic validation
was added to reason about relationships among addresses, collision domains,
gateways, node roles, and routing declarations before any output was written.

The next iterations were driven by increasingly complex routed laboratories.
Simple host--router scenarios required only interfaces and gateways, whereas
multi-area OSPF, RIP, and multi-AS BGP scenarios required protocol-specific
identifiers, areas, autonomous-system numbers, neighbours, advertised prefixes,
and redistribution rules. Rather than replacing the earlier configuration
files whenever the representation grew, the implementation retained compatible
forms where practical. This reduced disruption to already working laboratories
and made them useful as regression inputs, at the cost of some redundancy in
the accepted FRRouting syntax. The same pragmatic principle motivated escape
mechanisms such as \texttt{commands} and \texttt{frr.config}: they allowed the
project to express specialised behaviour without first modelling every Linux or
FRRouting command in the declarative language.

Execution of the generated laboratories then introduced a different class of
requirements. Inspecting files was not enough to determine whether packets were
actually following the intended path or whether an application was available.
Optional Wireshark integration was added as a diagnostic aid, lifecycle
operations were kept separate from generation, and connectivity testing was
progressively extended from basic reachability towards application-level
verification. During development, a reachable target could still fail a service
test when the application process was not running; this made the distinction
between network reachability and service availability explicit in the final
measurement design. Similarly, sequential client tests were adequate for
checking individual hosts but were unsuitable once several clients had to
observe the same service condition. The test runner was therefore changed to a
round-based concurrent model.

The experimental phase continued the same feedback loop. Earlier runs showed
that application measurements could be contaminated by a network that had not
yet reached a usable routing state. This motivated explicit readiness checks
before collecting the baseline. As the repository accumulated several
configurations and generated laboratories, utilities were also adapted to avoid
assuming that only one laboratory existed and to use explicit targets or node
names where appropriate. Finally, after functional behaviour was stable enough
to support the controlled DoS experiment, development moved to verification of
the translator itself. The stress and robustness tests in
Chapter~\ref{chap:robustness-security-scalability} exposed new trade-offs,
including duplicated validation work, permissive unknown YAML keys, incomplete
analysis of arbitrary static routes, and the scalability cost of collision-domain
overlap checks.

Figure~\ref{fig:development-lifecycle} summarises the lifecycle used throughout
the project. Verification is not represented as a terminal activity: an
observed problem can create a new or refined requirement and start another
iteration.

\begin{figure}[htbp]
\centering
\begin{tikzpicture}[
    node distance=8mm and 14mm,
    state/.style={draw, rounded corners, align=center, minimum width=3.5cm, minimum height=8mm, font=\small},
    arrow/.style={-{Latex[length=2mm]}, thick}
]
\node[state] (need) {Problem or\\new requirement};
\node[state, right=of need] (design) {Design choice\\and alternative analysis};
\node[state, below=of design] (impl) {Implementation\\or refactoring};
\node[state, left=of impl] (verify) {Laboratory execution\\and verification};
\node[state, below=of verify] (observe) {Problem, limitation,\\or new need observed};

\draw[arrow] (need) -- (design);
\draw[arrow] (design) -- (impl);
\draw[arrow] (impl) -- (verify);
\draw[arrow] (verify) -- (observe);
\draw[arrow] (observe.west) -| ++(-1.0,0) |- (need.west);
\end{tikzpicture}
\caption{Iterative software-development lifecycle followed by the project.}
\label{fig:development-lifecycle}
\end{figure}

The main iterations can be related directly to concrete features of the final
framework, as summarised in Table~\ref{tab:development-evolution}. The table is
not intended as a strict chronological log of every code change; instead, it
identifies the recurring engineering pattern that connected requirements,
observed problems, and implementation decisions.

{\small
\begin{longtable}{>{\raggedright\arraybackslash}p{0.18\textwidth}>{\raggedright\arraybackslash}p{0.34\textwidth}>{\raggedright\arraybackslash}p{0.38\textwidth}}
\caption{Main development iterations and resulting design choices.}
\label{tab:development-evolution}\\
\toprule
\textbf{Development stage} & \textbf{Need or problem observed} & \textbf{Resulting design or feature} \\
\midrule
\endfirsthead
\toprule
\textbf{Development stage} & \textbf{Need or problem observed} & \textbf{Resulting design or feature} \\
\midrule
\endhead
Initial prototype & Repetitive and duplicated information across several Kathar\'a artefacts. & A declarative YAML source from which standard Kathar\'a files are generated. \\
Validation & Generated files could remain syntactically plausible while describing an incoherent network. & Separate structural and semantic validation before filesystem modification. \\
Routing extension & Larger laboratories required OSPF, RIP, BGP, multi-area, and multi-AS information. & Protocol-aware validation and FRRouting generation, while retaining compatible declaration forms. \\
Extensibility & Some node behaviour could not be represented by the current declarative fields. & Generic \texttt{commands}, imported node resources, and the \texttt{frr.config} escape mechanism. \\
Runtime diagnosis & Static inspection did not show actual packet paths or all runtime failures. & Optional Wireshark integration and explicit post-start inspection with Linux and FRR tools. \\
Service verification & A reachable host or open TCP port did not necessarily mean that the application completed the requested operation. & Application-level request validation and distinct ICMP, TCP, and application metrics. \\
Concurrent testing & Sequential client execution observed the service at different times. & Round-based concurrent measurements across legitimate clients. \\
Experiment orchestration & Early measurements could begin before the routed topology was ready. & Readiness checks, isolated experiment outputs, and explicit baseline/adverse phases. \\
Final verification & Functional success did not reveal all robustness or scalability limits. & Stress testing and analysis of permissive or expensive implementation choices. \\
\bottomrule
\end{longtable}
}

\section{System Architecture and Repository Organisation}
\label{sec:system-architecture-repository}

The architecture is organised as a sequential pipeline with explicit
responsibility boundaries. Figure~\ref{fig:framework-architecture} shows the
main flow from source configuration to the running laboratory.

\begin{figure}[htbp]
\centering
\begin{tikzpicture}[
    node distance=0.55cm,
    box/.style={rectangle,draw,rounded corners,minimum width=11cm,minimum height=0.85cm,align=center}
]
\node[box] (source) {YAML source configuration and optional node-resource directories};
\node[box,below=of source] (loader) {Path resolution and safe YAML loading (\texttt{Lexer.py})};
\node[box,below=of loader] (structure) {Structural validation (\texttt{Parser.py})};
\node[box,below=of structure] (semantic) {Semantic validation and validation report (\texttt{Parser.py})};
\node[box,below=of semantic] (generation) {Kathar\'a and FRRouting generation (\texttt{gen\_lab.py})};
\node[box,below=of generation] (output) {Generated laboratory under \texttt{labs/<lab\_name>/}};
\node[box,below=of output] (lifecycle) {Start, inspect, clean, and remove through Kathar\'a and wrapper scripts};

\draw[-{Latex[length=3mm]}] (source) -- (loader);
\draw[-{Latex[length=3mm]}] (loader) -- (structure);
\draw[-{Latex[length=3mm]}] (structure) -- (semantic);
\draw[-{Latex[length=3mm]}] (semantic) -- (generation);
\draw[-{Latex[length=3mm]}] (generation) -- (output);
\draw[-{Latex[length=3mm]}] (output) -- (lifecycle);
\end{tikzpicture}
\caption{High-level architecture of the laboratory-generation framework.}
\label{fig:framework-architecture}
\end{figure}

The pipeline does not construct a persistent hierarchy of domain objects such
as \texttt{Node} or \texttt{Interface} classes. Introducing a typed internal
model was considered a possible way to make invariants more explicit, but it
would also have added another conversion layer between YAML and the relatively
compact output model required by the prototype. The implemented version
therefore keeps the validated configuration as a nested combination of Python
dictionaries and lists. This choice reduced implementation overhead and made
new YAML properties easy to prototype, but it also means that the validator is
responsible for establishing all assumptions later used by the generator. The
limitations of this choice become visible again in the discussion of strict
schemas and unknown keys in Chapter~\ref{chap:robustness-security-scalability}.

The root-level \texttt{generate\_lab.py} file is intentionally minimal: it
imports and invokes \texttt{main()} from
\texttt{src\_gen\_lab/gen\_lab.py}. Configuration discovery and
deserialisation are handled by \texttt{Lexer.py}. Despite its historical
filename, this component does not perform token-by-token lexical analysis; its
architectural role is to resolve the supplied configuration and load YAML.
\texttt{Parser.py} performs the minimum structural checks and the complete
semantic analysis. It defines the supported node and protocol sets, the
identifier pattern, the validation issue representation, and the report used
by both validation-only and normal generation modes. The generation module,
\texttt{gen\_lab.py}, then selects node images, generates \texttt{lab.conf},
writes startup files, creates FRRouting configuration, copies matching
resource directories, integrates Wireshark when requested, and reports how the
resulting laboratory can be started and stopped.

Lifecycle management is intentionally kept outside the generator itself.
\texttt{start.sh} verifies that the supplied laboratory directory exists,
changes to that directory, and executes \texttt{kathara lstart}.
\texttt{stop.sh} performs the corresponding \texttt{kathara lclean}, while
\texttt{del\_lab.sh} removes one generated directory or, when explicitly
requested, the contents of the laboratory output directory. This separation
keeps generation, runtime management, and filesystem deletion as distinct
operations.

The repository contains both the core framework and the later experimental
components. The following view includes only the files relevant to the
configuration, generation, and lifecycle architecture discussed in the first
part of the thesis.

\begin{lstlisting}[caption={Repository structure relevant to the framework core.},label={lst:repository-structure}]
Thesis/
|-- configs/
|   |-- Test1.yml
|   |-- Test_enterprise_multi_as_structured.yml
|   |-- client-router-server.yml
|   |-- dos_lab.yml
|   |-- mitigation.yml
|   |-- multi-area.yml
|   `-- multi_AS.yml
|-- src_gen_lab/
|   |-- Lexer.py
|   |-- Parser.py
|   |-- __init__.py
|   `-- gen_lab.py
|-- labs/
|-- generate_lab.py
|-- start.sh
|-- stop.sh
|-- del_lab.sh
`-- README.md
\end{lstlisting}

The \texttt{configs/} directory is the persistent source area. It contains YAML
topologies and may also contain directories whose names match declared nodes.
The \texttt{src\_gen\_lab/} package contains the implementation, whereas
\texttt{labs/} contains generated output and is intentionally not expanded in
the source-tree representation. This organisation expresses an important
architectural distinction: source configuration is maintained independently
from generated files and from the runtime state of the laboratory.

\section{Operational Flow and Design Principles}
\label{sec:operational-flow-design}

The control flow begins when the user invokes the command-line entry point.
The argument parser collects the requested configuration, the input and output
directories, overwrite behaviour, validation mode, resource-import option, and
Wireshark selection. The configuration path is then resolved. If the supplied
argument already exists as a path, it is used directly; otherwise, the loader
searches the configured directory for matching \texttt{.yml} and
\texttt{.yaml} files. The YAML content is deserialised into ordinary Python
objects.

Structural validation runs before semantic analysis so that subsequent
functions do not have to operate on a document that lacks mandatory elements
such as \texttt{nodes} or \texttt{interfaces}. Semantic validation then checks
the relationships inside the model and deliberately collects all detectable
errors and warnings instead of stopping at the first issue. If at least one
error exists, generation is not started. If validation succeeds and
\cliopt{validate-only} is active, the process terminates without modifying
the filesystem. Otherwise, the generator prepares the selected output
directory and writes all derived artefacts.

The generated directory is independent from the running state. Generation
creates files; starting the laboratory creates containers and collision
domains; stopping it removes the runtime environment; deleting the generated
directory is yet another separate operation. Keeping these stages distinct
prevents the project from conflating the declarative source, the generated
representation, and the executing instance.

Several design principles follow directly from this flow. First, the YAML
file is treated as the declarative source of truth for topology and common
node configuration. Generated files are derived output, which reduces
duplication and allows the laboratory to be recreated through clean
generation. Second, the framework deliberately generates standard Kathar\'a
files instead of controlling containers directly through the Docker API or a
private orchestration layer. Direct Docker control would have provided more
runtime control, but it would also have duplicated responsibilities already
implemented by Kathar\'a and tied the project to its own execution path. By
generating ordinary Kathar\'a artefacts, the framework remains an input and
validation layer rather than a replacement runtime; the output stays
compatible with the normal Kathar\'a workflow and can be inspected or executed
manually \cite{kathara_manpage}.

Validation is also deliberately placed before filesystem modification. A
structural or semantic error should not leave a partially generated
laboratory, and the validation-only option makes the same checks available for
manual review or automated workflows. Destructive operations follow the same
principle of explicitness: an existing laboratory directory is not silently
replaced. The user must select \cliopt{force} to write into existing output
or \cliopt{clean} to remove and recreate it, making the potentially-
destructive behaviour visible at the command line.

Resource import uses a convention rather than an additional mapping file. A
node-specific directory is associated with a node when their names match. This
choice is simple and readable, but it creates a deliberate naming convention;
when import checking is enabled, directories that do not match declared nodes
are reported and skipped. Diagnostics are likewise optional rather than
mandatory. Wireshark integration can be disabled non-interactively, while the
user may attach the monitoring node only to the collision domains required for
debugging or inspection. The core topology can therefore be generated without
monitoring, and observation can be added only when needed.

Together, these decisions support the non-functional objectives identified in
Section~\ref{sec:problem-scope-requirements}. Reproducibility is obtained by maintaining the YAML
source and regenerating output from it; transparency and compatibility are
preserved through ordinary Kathar\'a and FRRouting files; safety is improved by
validation-first and explicit overwrite semantics; modularity is reflected in
the separation among loading, validation, generation, monitoring selection,
and lifecycle operations; and traceability is provided by validation messages
that refer to the affected node, interface, collision domain, neighbour, or
field.


\section{LLM-Assisted Configuration and Deterministic Validation}
\label{sec:llm-role}

General-purpose LLMs can assist with configuration authoring, but authoring,
validation, and experiment execution are distinct responsibilities. The
relevant question is therefore not whether an LLM can emit configuration text,
but which guarantees are required before that text is used in a repeatable
experiment. In this thesis, RQ5 is addressed as a design and discussion
question through architectural analysis and comparison with the literature;
no empirical benchmark of LLM-generated configurations is claimed.

An LLM can translate informal intent into a candidate topology, allocate
addresses, draft startup commands, produce FRRouting fragments, or create the
YAML source accepted by the framework. Recent studies examine these
capabilities in network configuration and intent-based management
\cite{wang2024netconfeval,lira2025llmnetcfg,liu2025llmnetworking}. Such output
is nevertheless candidate input rather than evidence that all addresses,
gateways, routing relationships, node names, and cross-file references satisfy
the laboratory invariants.

The framework provides the complementary deterministic boundary. It applies
the same structural and semantic checks on every invocation, derives
coordinated Kathar\'a and FRRouting artefacts from one validated representation,
and connects the generated laboratory to runtime measurements. This separation
is important because technically plausible generated output may still contain
functional or security weaknesses and therefore requires explicit checking
\cite{pearce2022copilotsecurity}. The comparison in
Table~\ref{tab:llm-framework-roles} summarises these complementary roles.

\begin{table}[htbp]
\centering
\small
\begin{tabular}{p{0.23\textwidth} p{0.32\textwidth} p{0.32\textwidth}}
\toprule
\textbf{Activity} & \textbf{Typical LLM contribution} & \textbf{Framework contribution} \\
\midrule
Expressing a network requirement & Accepts natural-language intent and can ask or answer clarifying questions. & Requires an explicit YAML representation; it does not interpret unrestricted natural language. \\
Drafting topology and routing data & Can propose node roles, addresses, protocol parameters, and configuration fragments. & Provides the supported field model and turns accepted fields into standard Kathar\'a artefacts. \\
Checking known invariants & Can review a draft and suggest likely mistakes, but the result remains model-generated advice. & Applies explicit structural and semantic rules and blocks generation on recognised errors. \\
Maintaining cross-file consistency & Can generate several files, but consistency depends on the complete prompt, context, and generated output. & Derives multiple files from the same validated source representation. \\
Handling unusual or unmodelled behaviour & Can suggest arbitrary Linux or FRRouting commands and explain advanced features. & Supports controlled escape mechanisms such as \texttt{commands} and \texttt{frr.config}, but cannot fully validate their semantics. \\
Runtime diagnosis & Can interpret logs, command output, packet traces, and measured results supplied by the user. & Provides lifecycle operations, measurement tools, and optional packet capture that produce the evidence to be interpreted. \\
Reproducibility & Requires preservation of the selected prompt, context, model assumptions, and generated artefacts if the exact interaction is important. & Preserves a versionable source configuration and deterministically regenerates the supported laboratory structure. \\
\bottomrule
\end{tabular}
\caption{Complementary roles of LLM assistance and the developed framework.}
\label{tab:llm-framework-roles}
\end{table}

The comparison does not establish that either component is universally
superior. An LLM is suited to interpreting ambiguous intent, drafting candidate
configurations, explaining syntax, and helping a user interpret diagnostics.
The framework is narrower: it enforces the implemented configuration model,
maintains consistency among generated files, and produces the runtime evidence
required for evaluation.

\paragraph{Combined authoring and validation workflow.}
The workflow does not require a model embedded in the framework. A user may ask
an LLM to draft or refine YAML and then submit that candidate to the same
validation pipeline used for manually written input. Invalid candidates are
rejected with structured diagnostics; accepted candidates are translated into
standard Kathar\'a artefacts and executed. Logs, \texttt{vtysh} output,
measurements, and packet captures can subsequently support human or LLM-assisted
analysis.

\begin{figure}[htbp]
\centering
\begin{tikzpicture}[
    node distance=0.72cm,
    every node/.style={font=\small},
    box/.style={draw, rounded corners, minimum width=5.2cm, minimum height=0.65cm, align=center},
    optional/.style={box, dashed},
    arrow/.style={-{Latex[length=2mm]}, thick}
]
\node[box] (intent) {Experiment idea or natural-language network requirement};
\node[optional, below=of intent] (llm1) {Optional LLM assistance: propose or refine YAML};
\node[box, below=of llm1] (yaml) {Candidate declarative YAML configuration};
\node[box, below=of yaml] (validate) {Framework structural and semantic validation};
\node[box, below=of validate] (generate) {Deterministic Kathar\'a / FRRouting generation};
\node[box, below=of generate] (runtime) {Kathar\'a execution, runtime checks, and measurements};
\node[optional, below=of runtime] (llm2) {Optional human/LLM analysis of diagnostics and results};

\draw[arrow] (intent) -- (llm1);
\draw[arrow] (llm1) -- (yaml);
\draw[arrow] (yaml) -- (validate);
\draw[arrow] (validate) -- node[right]{\scriptsize valid} (generate);
\draw[arrow] (generate) -- (runtime);
\draw[arrow] (runtime) -- (llm2);
\draw[arrow] (validate.east) -- ++(1.05,0) node[right,align=left]{\scriptsize errors / warnings} |- (llm1.east);
\end{tikzpicture}
\caption{Recommended workflow in which LLM assistance remains optional and the framework provides the deterministic validation and generation boundary.}
\label{fig:llm-hybrid-workflow}
\end{figure}

The feedback path in Figure~\ref{fig:llm-hybrid-workflow} allows validation
diagnostics to guide a bounded human- or LLM-assisted repair step. The accepted
YAML remains the authoritative, versionable input, while measurements and logs
from the emulated network remain the evidence used to evaluate the experiment.

This architectural complementarity should not be interpreted as an empirical
claim about LLM performance or as proof of complete network correctness. The
current parser handles unknown keys permissively, and escape mechanisms such as
\texttt{commands} and \texttt{frr.config} may contain behaviour outside its
semantic model. Stricter schemas, repair policies, approval boundaries, and
runtime verification are therefore future work, discussed further in
\Cref{sec:future-llm-integration}.

\section{Architectural Boundaries and Chapter Summary}
\label{sec:architectural-boundaries-summary}

The architecture deliberately remains specific to Kathar\'a and its file
conventions. Supporting a different emulator would require a new output layer,
and the current implementation should not be interpreted as a general network
orchestration platform. The in-memory configuration is also based on nested
mappings rather than explicit domain classes. This keeps the code compact for
the present scope, but it makes disciplined validation and consistent key
handling essential. A future version could introduce a formal schema or
explicit classes for laboratories, nodes, interfaces, and routes.

The YAML model does not attempt to expose every Linux networking option or
every FRRouting command. Advanced configurations can still be expressed
through custom startup commands or by supplying a complete manual FRR
configuration through the supported escape hatch. This retains flexibility,
but semantic validation can only reason about the configuration properties it
understands explicitly. Generated files may also be edited manually, but such
changes do not become persistent source information: a clean regeneration will
remove them. Long-term modifications must therefore be made in the YAML source
or in the corresponding node-resource directory.

Within these boundaries, the framework implements a validation-first pipeline
that converts a declarative source configuration into standard Kathar\'a
artefacts while preserving a clear distinction among source files, generated
output, and running containers. The comparison with LLM-assisted configuration
adds another boundary: natural-language reasoning and draft generation may
occur outside the framework, but the accepted YAML source remains the explicit
interface to deterministic validation and generation. The architecture defines
responsibilities for configuration loading, structural and semantic validation,
laboratory generation, optional diagnostics, and lifecycle management, while
the functional and non-functional requirements establish the behaviour expected
from each stage.

The next chapter examines the configuration language and the validation
subsystem in detail. It describes the accepted YAML structures, path
resolution and loading, structural checks, semantic rules, and the validation
report used to decide whether laboratory generation may proceed.
